%==============================================================================
% Tento soubor použijte jako základ
% This file should be used as a base for the thesis
% Autoři / Authors: 2008 Michal Bidlo, 2022 Jaroslav Dytrych
% Kontakt pro dotazy a připomínky: sablona@fit.vutbr.cz
% Contact for questions and comments: sablona@fit.vutbr.cz
%==============================================================================
% kódování: UTF-8 (zmena prikazem iconv, recode nebo cstocs)
% encoding: UTF-8 (you can change it by command iconv, recode or cstocs)
%------------------------------------------------------------------------------
% zpracování / processing: make, make pdf, make clean
%==============================================================================
% Soubory, které je nutné upravit nebo smazat: / Files which have to be edited or deleted:
%   projekt-20-literatura-bibliography.bib - literatura / bibliography
%   projekt-01-kapitoly-chapters.tex - obsah práce / the thesis content
%   projekt-01-kapitoly-chapters-en.tex - obsah práce v angličtině / the thesis content in English
%   projekt-30-prilohy-appendices.tex - přílohy / appendices
%   projekt-30-prilohy-appendices-en.tex - přílohy v angličtině / appendices in English
%==============================================================================
\documentclass[]{fitthesis} % bez zadání - pro začátek práce, aby nebyl problém s překladem
%\documentclass[english]{fitthesis} % without assignment - for the work start to avoid compilation problem
%\documentclass[zadani]{fitthesis} % odevzdani do IS VUT a/nebo tisk s barevnými odkazy - odkazy jsou barevné
%\documentclass[english,zadani]{fitthesis} % for submission to the IS VUT and/or print with color links - links are color
%\documentclass[zadani,print]{fitthesis} % pro černobílý tisk - odkazy jsou černé
%\documentclass[english,zadani,print]{fitthesis} % for the black and white print - links are black
%\documentclass[zadani,cprint]{fitthesis} % pro barevný tisk - odkazy jsou černé, znak VUT barevný
%\documentclass[english,zadani,cprint]{fitthesis} % for the print - links are black, logo is color
% * Je-li práce psaná v anglickém jazyce, je zapotřebí u třídy použít 
%   parametr english následovně:
%   If thesis is written in English, it is necessary to use 
%   parameter english as follows:
%      \documentclass[english]{fitthesis}
% * Je-li práce psaná ve slovenském jazyce, je zapotřebí u třídy použít 
%   parametr slovak následovně:
%   If the work is written in the Slovak language, it is necessary 
%   to use parameter slovak as follows:
%      \documentclass[slovak]{fitthesis}
% * Je-li práce psaná v anglickém jazyce se slovenským abstraktem apod., 
%   je zapotřebí u třídy použít parametry english a enslovak následovně:
%   If the work is written in English with the Slovak abstract, etc., 
%   it is necessary to use parameters english and enslovak as follows:
%      \documentclass[english,enslovak]{fitthesis}

% Základní balíčky jsou dole v souboru šablony fitthesis.cls
% Basic packages are at the bottom of template file fitthesis.cls
% zde můžeme vložit vlastní balíčky / you can place own packages here


% Pro seznam zkratek lze využít balíček Glossaries - nutno odkomentovat i níže a při kompilaci z konzoly i v Makefile (plnou verzi pro Perl, nebo lite)
% The Glossaries package can be used for the list of abbreviations - it is necessary to uncomment also below. When compiling from the console also in the Makefile (full version for Perl or lite)
%\usepackage{glossaries}
%\usepackage{glossary-superragged}
%\makeglossaries 

% Nastavení cesty k obrázkům
% Setting of a path to the pictures
%\graphicspath{{obrazky-figures/}{./obrazky-figures/}}
%\graphicspath{{obrazky-figures/}{../obrazky-figures/}}

%---rm---------------
\renewcommand{\rmdefault}{lmr}%zavede Latin Modern Roman jako rm / set Latin Modern Roman as rm
%---sf---------------
\renewcommand{\sfdefault}{qhv}%zavede TeX Gyre Heros jako sf
%---tt------------
\renewcommand{\ttdefault}{lmtt}% zavede Latin Modern tt jako tt

% vypne funkci šablony, která automaticky nahrazuje uvozovky,
% aby nebyly prováděny nevhodné náhrady v popisech API apod.
% disables function of the template which replaces quotation marks
% to avoid unnecessary replacements in the API descriptions etc.
\csdoublequotesoff

\usepackage{url}

% =======================================================================
% balíček "hyperref" vytváří klikací odkazy v pdf, pokud tedy použijeme pdflatex
% problém je, že balíček hyperref musí být uveden jako poslední, takže nemůže
% být v šabloně
% "hyperref" package create clickable links in pdf if you are using pdflatex.
% Problem is that this package have to be introduced as the last one so it 
% can not be placed in the template file.
\ifWis
\ifx\pdfoutput\undefined % nejedeme pod pdflatexem / we are not using pdflatex
\else
  \usepackage{color}
  \usepackage[unicode,colorlinks,hyperindex,plainpages=false,pdftex]{hyperref}
  \definecolor{hrcolor-ref}{RGB}{223,52,30}
  \definecolor{hrcolor-cite}{HTML}{2F8F00}
  \definecolor{hrcolor-urls}{HTML}{092EAB}
  \hypersetup{
	linkcolor=hrcolor-ref,
	citecolor=hrcolor-cite,
	filecolor=magenta,
	urlcolor=hrcolor-urls
  }
  \def\pdfBorderAttrs{/Border [0 0 0] }  % bez okrajů kolem odkazů / without margins around links
  \pdfcompresslevel=9
\fi
\else % pro tisk budou odkazy, na které se dá klikat, černé / for the print clickable links will be black
\ifx\pdfoutput\undefined % nejedeme pod pdflatexem / we are not using pdflatex
\else
  \usepackage{color}
  \usepackage[unicode,colorlinks,hyperindex,plainpages=false,pdftex,urlcolor=black,linkcolor=black,citecolor=black]{hyperref}
  \definecolor{links}{rgb}{0,0,0}
  \definecolor{anchors}{rgb}{0,0,0}
  \def\AnchorColor{anchors}
  \def\LinkColor{links}
  \def\pdfBorderAttrs{/Border [0 0 0] } % bez okrajů kolem odkazů / without margins around links
  \pdfcompresslevel=9
\fi
\fi
% Řešení problému, kdy klikací odkazy na obrázky vedou za obrázek
% This solves the problems with links which leads after the picture
\usepackage[all]{hypcap}


% Informace o práci/projektu / Information about the thesis
%---------------------------------------------------------------------------
\projectinfo{
  %Prace / Thesis
  project={BP},            %typ práce BP/SP/DP/DR  / thesis type (SP = term project)
  year={2025},             % rok odevzdání / year of submission
  date=\today,             % datum odevzdání / submission date
  %Nazev prace / thesis title
  title.cs={Název práce},  % název práce v češtině či slovenštině (dle zadání) / thesis title in czech language (according to assignment)
  title.en={Thesis title}, % název práce v angličtině / thesis title in english
  %title.length={14.5cm}, % nastavení délky bloku s titulkem pro úpravu zalomení řádku (lze definovat zde nebo níže) / setting the length of a block with a thesis title for adjusting a line break (can be defined here or below)
  %sectitle.length={14.5cm}, % nastavení délky bloku s druhým titulkem pro úpravu zalomení řádku (lze definovat zde nebo níže) / setting the length of a block with a second thesis title for adjusting a line break (can be defined here or below)
  %dectitle.length={14.5cm}, % nastavení délky bloku s titulkem nad prohlášením pro úpravu zalomení řádku (lze definovat zde nebo níže) / setting the length of a block with a thesis title above declaration for adjusting a line break (can be defined here or below)
  %Autor / Author
  author.name={Jméno},   % jméno autora / author name
  author.surname={Příjmení},   % příjmení autora / author surname 
  %author.title.p={Bc.}, % titul před jménem (nepovinné) / title before the name (optional)
  %author.title.a={Ph.D.}, % titul za jménem (nepovinné) / title after the name (optional)
  %Ustav / Department
  department={UPGM}, % doplňte příslušnou zkratku dle ústavu na zadání: UPSY/UIFS/UITS/UPGM / fill in appropriate abbreviation of the department according to assignment: UPSY/UIFS/UITS/UPGM
  % Školitel / supervisor
  supervisor.name={Jméno},   % jméno školitele / supervisor name 
  supervisor.surname={Příjmení},   % příjmení školitele / supervisor surname
  supervisor.title.p={prof. RNDr.},   %titul před jménem (nepovinné) / title before the name (optional)
  supervisor.title.a={Ph.D.},    %titul za jménem (nepovinné) / title after the name (optional)
  % Klíčová slova / keywords
  keywords.cs={Sem budou zapsána jednotlivá klíčová slova v českém (slovenském) jazyce, oddělená čárkami.}, % klíčová slova v českém či slovenském jazyce / keywords in czech or slovak language
  keywords.en={Sem budou zapsána jednotlivá klíčová slova v anglickém jazyce, oddělená čárkami.}, % klíčová slova v anglickém jazyce / keywords in english
  %keywords.en={Here, individual keywords separated by commas will be written in English.},
  % Abstrakt / Abstract
  abstract.cs={Do tohoto odstavce bude zapsán výtah (abstrakt) práce v českém (slovenském) jazyce.}, % abstrakt v českém či slovenském jazyce / abstract in czech or slovak language
  abstract.en={Do tohoto odstavce bude zapsán výtah (abstrakt) práce v anglickém jazyce.}, % abstrakt v anglickém jazyce / abstract in english
  %abstract.en={An abstract of the work in English will be written in this paragraph.},
  % Prohlášení (u anglicky psané práce anglicky, u slovensky psané práce slovensky; u projektové praxe lze zakomentovat) / Declaration (for thesis in english should be in english; for project practice can be commented out)
  declaration={Prohlašuji, že jsem tuto bakalářskou práci vypracoval samostatně pod vedením pana X...
Další informace mi poskytli...
Uvedl jsem všechny literární prameny, publikace a další zdroje, ze kterých jsem čerpal.},
  %declaration={I hereby declare that this Bachelor's thesis was prepared as an original work by the author under the supervision of Mr. X
% The supplementary information was provided by Mr. Y
% I have listed all the literary sources, publications and other sources, which were used during the preparation of this thesis.},
  % Poděkování (nepovinné, nejlépe v jazyce práce; nechcete-li, zakomentujte pro skrytí nadpisu) / Acknowledgement (optional, ideally in the language of the thesis; comment out for hiding including heading)
  acknowledgment={V této sekci je možno uvést poděkování vedoucímu práce a těm, kteří poskytli odbornou pomoc
(externí zadavatel, konzultant apod.).},
  %acknowledgment={Here it is possible to express thanks to the supervisor and to the people which provided professional help
%(external submitter, consultant, etc.).},
  % Rozšířený abstrakt (cca 3 normostrany) - lze definovat zde nebo níže / Extended abstract (approximately 3 standard pages) - can be defined here or below
  %extendedabstract={Do tohoto odstavce bude zapsán rozšířený výtah (abstrakt) práce v českém (slovenském) jazyce.},
  %extabstract.odd={true}, % Začít rozšířený abstrakt na liché stránce? / Should extended abstract start on the odd page?
  %faculty={FIT}, % FIT/FEKT/FSI/FA/FCH/FP/FAST/FAVU/USI/DEF
  faculty.cs={Fakulta informačních technologií}, % Fakulta v češtině - pro využití této položky výše zvolte fakultu DEF / Faculty in Czech - for use of this entry select DEF above
  faculty.en={Faculty of Information Technology}, % Fakulta v angličtině - pro využití této položky výše zvolte fakultu DEF / Faculty in English - for use of this entry select DEF above
  department.cs={Ústav matematiky}, % Ústav v češtině - pro využití této položky výše zvolte ústav DEF nebo jej zakomentujte / Department in Czech - for use of this entry select DEF above or comment it out
  department.en={Institute of Mathematics} % Ústav v angličtině - pro využití této položky výše zvolte ústav DEF nebo jej zakomentujte / Department in English - for use of this entry select DEF above or comment it out
}

% Rozšířený abstrakt (cca 3 normostrany) - lze definovat zde nebo výše / Extended abstract (approximately 3 standard pages) - can be defined here or above
%\extendedabstract{Do tohoto odstavce bude zapsán výtah (abstrakt) práce v českém (slovenském) jazyce.}
% Začít rozšířený abstrakt na liché stránce? / Should extended abstract start on the odd page?
%\extabstractodd{true}

% nastavení délky bloku s titulkem pro úpravu zalomení řádku - lze definovat zde nebo výše / setting the length of a block with a thesis title for adjusting a line break - can be defined here or above
%\titlelength{14.5cm}
% nastavení délky bloku s druhým titulkem pro úpravu zalomení řádku - lze definovat zde nebo výše / setting the length of a block with a second thesis title for adjusting a line break - can be defined here or above
%\sectitlelength{14.5cm}
% nastavení délky bloku s titulkem nad prohlášením pro úpravu zalomení řádku - lze definovat zde nebo výše / setting the length of a block with a thesis title above declaration for adjusting a line break - can be defined here or above
%\dectitlelength{14.5cm}

% řeší první/poslední řádek odstavce na předchozí/následující stránce
% solves first/last row of the paragraph on the previous/next page
\clubpenalty=10000
\widowpenalty=10000

% checklist
\newlist{checklist}{itemize}{1}
\setlist[checklist]{label=$\square$}

% Kompilace po částech (rychlejší, ale v náhledu nemusí být vše aktuální)
% Compilation piecewise (faster, but not all parts in preview will be up-to-date)
% Další informace viz / For more information see https://www.overleaf.com/learn/latex/Multi-file_LaTeX_projects
% \usepackage{subfiles}

% Nechcete-li, aby se u oboustranného tisku roztahovaly mezery pro zaplnění stránky, odkomentujte následující řádek / If you do not want enlarged spacing for filling of the pages in case of duplex printing, uncomment the following line
% \raggedbottom

\begin{document}
  % Vysazeni titulnich stran / Typesetting of the title pages
  % ----------------------------------------------
  \maketitle
  % Obsah
  % ----------------------------------------------
  \setlength{\parskip}{0pt}

  {\hypersetup{hidelinks}\tableofcontents}
  
  % Seznam obrazku a tabulek (pokud prace obsahuje velke mnozstvi obrazku, tak se to hodi)
  % List of figures and list of tables (if the thesis contains a lot of pictures, it is good)
  \ifczech
    \renewcommand\listfigurename{Seznam obrázků}
  \fi
  \ifslovak
    \renewcommand\listfigurename{Zoznam obrázkov}
  \fi
  {\hypersetup{hidelinks}\listoffigures}
  
  \ifczech
    \renewcommand\listtablename{Seznam tabulek}
  \fi
  \ifslovak
    \renewcommand\listtablename{Zoznam tabuliek}
  \fi
  % {\hypersetup{hidelinks}\listoftables}

  % Seznam zkratek / List of abbreviations
  %\ifczech
  %  \renewcommand*\glossaryname{Seznam zkratek}%
  %  \renewcommand*\entryname{Zkratka}
  %  \renewcommand*\descriptionname{Význam}
  %\fi
  %\ifslovak
  %  \renewcommand*\glossaryname{Zoznam skratiek}%
  %  \renewcommand*\entryname{Skratka}
  %  \renewcommand*\descriptionname{Význam}
  %\fi
  %\ifenglish
  %  \renewcommand*\glossaryname{List of abbreviations}%
  %  \renewcommand*\entryname{Abbreviation}
  %  \renewcommand*\descriptionname{Meaning}
  %\fi
  % Definice zkratek - z textu se odkazují např. \Gls{TF–IDF}
  % Definition of abbreviations - referred from the text e.g. \Gls{TF–IDF}
  %\newglossaryentry{TF–IDF}
  %{
  %  name={TF–IDF},
  %  description={Term Frequency-Inverse Document Frequency}
  %}
  % 
  %\setglossarystyle{superragged}
  %\printglossaries


  \ifODSAZ
    \setlength{\parskip}{0.5\bigskipamount}
  \else
    \setlength{\parskip}{0pt}
  \fi

  % vynechani stranky v oboustrannem rezimu
  % Skip the page in the two-sided mode
  \iftwoside
    \cleardoublepage
  \fi

  % Text prace / Thesis text
  % ----------------------------------------------
  \ifenglish
    \input{projekt-01-kapitoly-chapters-en}
  \else
    % Tento soubor nahraďte vlastním souborem s obsahem práce.
%=========================================================================
% Autoři: Michal Bidlo, Bohuslav Křena, Jaroslav Dytrych, Petr Veigend a Adam Herout 2019

% Pro kompilaci po částech (viz projekt.tex), nutno odkomentovat a upravit
%\documentclass[../projekt.tex]{subfiles}
%\begin{document}

\chapter{Úvod}
Podpis je jednou z~nejstarších metod používanou pro ověření totožnosti, určení autorství či udělení právního souhlasu.
V~této práci půjde především o~ověření totožnosti, které spadá pod obor zvaný biometrie.
Biometrie se zabývá analýzou biologických a behaviorálních charakteristik používaných mimo jiné k~autentizaci. 
Zkoumáním pravosti podpisu a jiných textů se zabývá písmoznalectví neboli grafognózie.
U~podpisu lze analyzovat statické a dynamické parametry písma. 

%Čím se tedy zabývá tato práce?
Podpis je v~praxi stále nejpoužívanější způsob autentizace. 
Podaří-li se důvěryhodně podpis napodobit, bude možnost obejít bezpečnostní zabezpečení. 
To představuje velké bezpečnostní riziko a nutnost upravení dosavadních bezpečnostních metod pro autentizaci. 
Bavíme se zde spíše o~digitálním podpisu. 
Na rozdíl od podpisu na papíře digitální podpis obsahuje i dynamické charakteristiky.
Bude tedy o něco složitější jej zfalšovat.

%TOLE JE SPIS ABSTRAKT
%Bude třeba sestavit snímací zařízení pro získaní potřebných parametrů podpisu.
%Ze statických parametrů budeme snímat tvar i tloušťku čar a celkový vzhled podpisu. 
%Z dynamických parametrů budeme snímat tlak, sklon a průběh tahů. 
%Získané data je potřeba vhodným způsobem zpracovat a provést potřebné výpočty. 
%Upravená data je poté potřeba převést do instrukcí pro mechanickou paži, která je použita k replikaci podpisu. 
%Výsledný podpis bude porovnáván s vlastnoručním podpisem. Cílem je co nejvíce se přiblížit vlastnoručnímu podpisu.

%\section{Představení tématu}
%\section{Cíle práce}

%\section{Struktura práce}

\chapter{Teoretické základy}
Pro plné pochopení problematiky je potřeba představit základy z~několika odvětví. 
V~této části budou shrnuty všechny potřebné informace. 

\section{Pojmy}
\label{sec:pojmy}
\begin{itemize}
  \item{\textbf{Biologické autentizační metody} - zkoumají biologické charakteristiky člověka, s~jimiž se narodí. Tyto charakteristiky má každý člověk unikátní a jsou neměnitelné. Patří mezi ně například otisky prstů, skeny sítnice či duhovky.}  
  \item{\textbf{Behaviorální autentizační metody} - porovnávají behaviorální charakteristiky související s~chováním člověka. Lze například porovnávat, jakým způsobem jedinec mluví nebo jak pohybuje očima.}
  \item{\textbf{Statické parametry písma} - jde o~parametry, které nezahrnují informace o~procesu psaní podpisu. Patří mezi ně například tvar a vzhled písma, umístění podpisu na stránce nebo také tloušťka čar.}
  \item{\textbf{Dynamické parametry písma} - tyto parametry se vytahují k~samotnému průběhu psaní podpisu. Například rychlost psaní, akcelerace, tlak nebo průběh tahů.}
\end{itemize}

\section{Definice podpisu a jeho charakteristiky}
Definice podpisu je z pohledu zákona poněkud problematická. 
Podpis jako takový není v~zákoně nijak přímo definován. 
Obecně je ale bráno za podpis vlastnoruční uvedení jména a příjmení, nebo jiného jedinečného a nezaměnitelného označení. %https://www.fulsoft.cz/33/komentar-zakona-89-2012-sb-obcansky-zakonik-561-uniqueidmRRWSbk196FNf8-jVUh4EtuvCojfP1Dmx8t50X_QR_yU_dtwjPItMA/

Charakteristiky podpisu je nepřeberné množství informací, podle kterých lze jednotlivé podpisy od sebe odlišit. 
Charakteristiky se dělí na statické a dynamické~\ref{sec:pojmy}. 
Jsou ovlivňovány jak fyzickým, tak i psychologickým stavem člověka v době podpisu.
Mezi charakteristiky patří například celkový vzhled podpisu, tloušťka čar, rychlost psaní, tlak, průběh a pořadí tahů a mnoho dalších, které ale pro nás v rámci této práce nebudou důležité. %TODO vymyslel jsem si lol

\section{Biometrická autentizace}
Biometrická autentizace probíhá na základě biologických a behaviorálních charakteristik člověka\ref{sec:pojmy}. 
Obě tyto skupiny znaků jsou pro každého člověka unikátní a nezaměnitelné. To je důvodem, proč je lze využít k autentizaci.
Existuje spousta biometrických metod pro autentizaci, kupříkladu otisk prstu, rozpoznání obličeje, sken duhovky či sítnice, nebo dokonce pomocí používání myši. 
Každá tato metoda má své výhody a nevýhody, zejména se hodnotí přesnost, cena, komfort používání, stálost a velikost vzorku.
Tato práce se zaměřuje na autentizaci podpisem.

\section{Výběr hardwarových komponent}

\section{Parametry podpisu}
V rámci této práce budou snímány následující parametry:
\begin{itemize}
  \item{\textbf{Celkový vzhled podpisu}}
  \item{\textbf{Pohyb pera} ze kterého bude vypočítán celkový průběh podpisu, především pak \textbf{pozice pera} a \textbf{rychlost psaní čar}.} %při psaní pomocí akcelerometru zabudovaném v MPU-6050 
  \item{\textbf{Sklon pera}} %pomocí gyroskopu, který je také součástí MPU-6050.}
  \item{\textbf{Tlak}} %pomocí čtveřice tlakových senzorů Interlink Electronics FSR® 400.}
\end{itemize}
Kombinací těchto nasnímaných parametrů by bylo možné dostatečně věrohodně napodobit původní vlastnoruční podpis. 

\begin{figure}[h]
\centering
\begin{minipage}{0.45\textwidth}
    \centering
    \includegraphics[width=\textwidth]{obrazky-figures/placeholder.pdf}
    \caption{Vzhled dynamického podpisu.}
    \label{fig:first-image}
\end{minipage}\hfill
\begin{minipage}{0.45\textwidth}
    \centering
    \includegraphics[width=\textwidth]{obrazky-figures/placeholder.pdf}
    \caption{Graf dynamických parametrů podpisu v čase.}
    \label{fig:second-image}
\end{minipage}
\end{figure}


\section{Rozpoznávání falzifikátů}
Rozpoznávání falzifikátů je u digitálního podpisu automatizované.
Probíhá na základě porovnávání uloženého vzorku v databázi s podpisem, kterým se daný člověk pokouší autentizovat. %použito správně slovo autentizce?
Referenční podpis je zprůměrován z několika vzorků, aby došlo k potlačení náhodných jevů. % zav_prace 17183 2.3
Tyto vzorky jsou poskytnuté danou osobou při registraci do systému.

Porovnávací algoritmus extrahuje důležité parametry podpisu, které poté porovnává s referenčním vzorem. 
Výsledkem takového porovnávání je určité procento shody parametrů.
Následně je vypočteno celkové procento shody podpisů a na jeho základě je autentizace úspěšná či nikoli.
Je velmi důležité, aby algoritmus měl určenou správnou procentuální hodnotu shody.
Pokud by byla špatně nastavena, bylo by příliš mnoho neúspěšných autentizací, které by měly být úspěšné, nebo naopak mnoho úspěšných autentizací, které měly být neúspěšné.


\begin{figure}[h]
\centering
\includegraphics[width=0.5\textwidth]{obrazky-figures/placeholder.pdf}
\caption{Průběh autentizace}
\label{fig:my-pdf}
\end{figure}

\subsection{rozpoznání statických falzifikátů}
Verifikace v off-line systému se skládá ze tří etap – předzpracování, extrakce biometrických
charakteristik a vyhodnocení. Během předzpracování jsou provedeny některé algoritmy, jako
jsou vyhlazování, zjednodušení, či normalizace. Extrakce probíhá různými způsoby, jedním
z nich je zjištění geometrických a topologických rysů – uzavřené smyčky, speciální body
nebo místa křížení (Obr. 15). Výsledné vyhodnocení je založeno na vyhodnocování vektorů
charakteristik, může být provedeno porovnáním významových bodů, klasifikátorem sousedů
nebo neuronovou sítí. (1, kap. 13.7) %https://theses.cz/id/odwy80/BP_Vchov.pdf (1) https://aleph.lib.vutbr.cz/F/QR556J3JR4JCE18FMRXDUGQ8FXVA2VL3NX3CYI9PYU2PRQCPTX-02400?func=item-global&doc_library=BUT01&doc_number=000088231&year=&volume=&sub_library=BXFIT

\subsection{rozpoznání dynamických falzifikátů}
V případě on-line verifikace jsou pomocí HW získány kromě statických i dynamické
charakteristiky – rychlost psaní, tlak pera, pořadí psaní jednotlivých částí apod. Již v průběhu
psaní podpisu jsou získávány informace o průměrné a maximální rychlosti psaní, měření
vlastností zakřivení tahů, poměr dlouhých a krátkých tahů, různé délky segmentů podpisu aj..
V každém časovém okamžiku vzniku podpisu jsou popisovány časovou funkcí další
dynamické charakteristiky – souřadnicové pozice hrotu psacího nástroje, rychlost (v),
zrychlení (a) a tlak hrotu na podložku (p) (Obr. 16). Následně jsou získané charakteristiky
porovnávány se šablonou uloženou v databázi, získá se tak míra ztotožnění. (1, kap. 13.8) %https://theses.cz/id/odwy80/BP_Vchov.pdf (1) https://aleph.lib.vutbr.cz/F/QR556J3JR4JCE18FMRXDUGQ8FXVA2VL3NX3CYI9PYU2PRQCPTX-02400?func=item-global&doc_library=BUT01&doc_number=000088231&year=&volume=&sub_library=BXFIT


\subsection{Normy identifikace pomocí podpisu}
Identifikace pomocí podpisu podléhá normám ISO/IEC 19794-7:2007 Signature/sign time
series data, ISO/IEC 19794-11:2013 Signature/sign processed dynamic data a ANSI/INCITS
395-2005 Signature/Sign Data.
ISO/IEC 19794-7:2007 popisuje dva formáty výměny dat pro podpis a písmo. První je obecný
formát, druhý je kompaktní, který lze využít pro smart karty. (48) % https://theses.cz/id/odwy80/BP_Vchov.pdf (48) https://www.iso.org/standard/38751.html
ISO/IEC 19794-11:2013 pro účely biometrického porovnání specifikuje obecný formát
výměny dat pro zpracované podpisové/znakové údaje o chování získané z časové řady
zachycené pomocí zařízení. (49) % https://theses.cz/id/odwy80/BP_Vchov.pdf (49) https://www.iso.org/standard/51824.html
ANSI/INCITS 395-2005 specifikuje formát výměny dat pro digitalizované písmo nebo podpis
pro účely biometrického zápisu, verifikace nebo identifikace. Tento formát je obecný a může
být využit v široké škále aplikací využívajících elektronický podpis. (50) % https://theses.cz/id/odwy80/BP_Vchov.pdf (50) https://webstore.ansi.org/standards/incits/ansiincits3952005

\chapter{Analýza současných řešení}


\section{Současná řešení v~oblasti biometrie}
\section{Systém, zařízení a kde se využívají, a jak jsou důvěryhodné}
\section{Hodnocení výhod a nevýhod existujících řešení}

\chapter{Návrh snímacího pera}
\section{Výběr senzorů a technologií}
\*subsection{EPS32}
\*subsection{MPU-6050 (akcelerometr a gyroskop)}
\*subsection{Tlakový senzor Interlink Electronics FSR® 400}
\section{Schéma a popis návrhu}
\section{Možnosti replikace podpisu}

%\chapter{Implementace prototypu}
%\section{Vývoj hardwaru}
%\section{Sbírání dat v~reálném čase}
%\section{Ukládání dat pro následnou replikaci podpisu}

%\chapter{Replikace podpisu}
%\section{Využití robotické ruky}
%\section{Transformace nasnímaných dat na kód pro replikaci}
%\section{Testování a výsledky replikace}

%\chapter{Hodnocení výsledků}
%\section{Analýza dosažených výsledků}
%\section{Porovnání s~očekáváními}
%\section{Diskuze o~spolehlivosti a přesnosti replikace}

%\chapter{Závěr}
%\section{Shrnutí hlavních poznatků}
%\section{Zhodnocení významu práce}
%\section{Budoucí perspektivy}

%\section{Možná vylepšení a rozšíření}
%\*subsection{Návrhy na zlepšení prototypu}
%\*subsection{Možnosti dalšího výzkumu a vývoje}


%===============================================================================

% Pro kompilaci po částech (viz projekt.tex) nutno odkomentovat
%\end{document}

  \fi
  
  % Kompilace po částech (viz výše, nutno odkomentovat a zakomentovat input výše)
  % Compilation piecewise (see above, it is necessary to uncomment it and comment out input above)
  %\subfile{chapters/projekt-01-uvod-introduction}
  % ...
  %\subfile{chapters/projekt-05-zaver-conclusion}

  % Pouzita literatura / Bibliography
  % ----------------------------------------------
\ifslovak
  \makeatletter
  \def\@openbib@code{\addcontentsline{toc}{chapter}{Literatúra}}
  \makeatother
  \bibliographystyle{bib-styles/Pysny/skplain}
\else
  \ifczech
    \makeatletter
    \def\@openbib@code{\addcontentsline{toc}{chapter}{Literatura}}
    \makeatother
    \bibliographystyle{bib-styles/Pysny/czplain}
  \else 
    \makeatletter
    \def\@openbib@code{\addcontentsline{toc}{chapter}{Bibliography}}
    \makeatother
    \bibliographystyle{bib-styles/Pysny/enplain}
  %  \bibliographystyle{alpha}
  \fi
\fi
  \begin{flushleft}
  \bibliography{projekt-20-literatura-bibliography}
  \end{flushleft}

  % vynechani stranky v oboustrannem rezimu
  % Skip the page in the two-sided mode
  \iftwoside
    \cleardoublepage
  \fi

  % Prilohy / Appendices
  % ---------------------------------------------
  \appendix
\ifczech
  \renewcommand{\appendixpagename}{Přílohy}
  \renewcommand{\appendixtocname}{Přílohy}
  \renewcommand{\appendixname}{Příloha}
\fi
\ifslovak
  \renewcommand{\appendixpagename}{Prílohy}
  \renewcommand{\appendixtocname}{Prílohy}
  \renewcommand{\appendixname}{Príloha}
\fi
%  \appendixpage

% vynechani stranky v oboustrannem rezimu
% Skip the page in the two-sided mode
%\iftwoside
%  \cleardoublepage
%\fi
  
\ifslovak
%  \section*{Zoznam príloh}
%  \addcontentsline{toc}{section}{Zoznam príloh}
\else
  \ifczech
%    \section*{Seznam příloh}
%    \addcontentsline{toc}{section}{Seznam příloh}
  \else
%    \section*{List of Appendices}
%    \addcontentsline{toc}{section}{List of Appendices}
  \fi
\fi
  \startcontents[chapters]
  \setlength{\parskip}{0pt} 
  % seznam příloh / list of appendices
  % \printcontents[chapters]{l}{0}{\setcounter{tocdepth}{2}}
  
  \ifODSAZ
    \setlength{\parskip}{0.5\bigskipamount}
  \else
    \setlength{\parskip}{0pt}
  \fi
  
  % vynechani stranky v oboustrannem rezimu
  \iftwoside
    \cleardoublepage
  \fi
  
  % Přílohy / Appendices
  \ifenglish
    \input{projekt-30-prilohy-appendices-en}
  \else
    \input{projekt-30-prilohy-appendices}
  \fi
  
  % Kompilace po částech (viz výše, nutno odkomentovat)
  % Compilation piecewise (see above, it is necessary to uncomment it)
  %\subfile{projekt-30-prilohy-appendices}
  
\end{document}
